%% ============================================================
%%  Lecture 6 -- Thursday, 11 June 2026
%%  Subfile of main.tex: compiles standalone OR as part of the book.
%% ============================================================
\documentclass[main.tex]{subfiles}

\begin{document}

\beginlecture{6}{Thursday, 11 June 2026}%
  {Monotone Convergence, the Squeeze Theorem, and the Upper and Lower Limits}

\epigraph{Of the small there is no smallest, but always a smaller; and of the large
there is always a larger.}{Anaxagoras, fragment B3 (5th c.\ bce)}

\begin{lecturecontents}{Contents of this lecture}
  \item \hyperlink{6.1}{6.1\quad Review of Completeness and Cauchy Sequences}
  \item \hyperlink{6.2}{6.2\quad Monotone sequences}
  \item \hyperlink{6.3}{6.3\quad Comparison and squeeze}
  \item \hyperlink{6.4}{6.4\quad The upper and lower limits}
  \item \hyperlink{6.5}{6.5\quad Some standard limits}
  \item \hyperlink{6.6}{6.6\quad Introduction to series}
\end{lecturecontents}

\bigskip
\noindent\textbf{Overview.} Picking up from completeness, this lecture assembles the basic
machinery for limits of real sequences. Monotonicity plus boundedness already forces
convergence; the comparison and squeeze principles let limits be read off from neighbours;
and for sequences that oscillate, the upper and lower limits $\limsup$ and $\liminf$ record
the band the sequence is eventually trapped in, coinciding exactly when the sequence
converges. After recording four limits that recur throughout the subject, we open the theory
of series: convergence through partial sums, the Cauchy criterion, and the first divergence test.
The companion text is Rudin, Chapter~3, especially \R{3.12}--\R{3.40}.

\clearpage
% ============================================================
\sub{6.1}{Review of Completeness and Cauchy Sequences}

\begin{Obj}{Definition}{6.1.1}{Complete metric space}
A metric space $(X,d)$ is \emph{complete} (\R{3.12}) if every Cauchy sequence in $X$
converges to a point of $X$.
\end{Obj}

\begin{ObjL}{Example}{6.1.2}{Complete and incomplete spaces}
\begin{itemize}[leftmargin=1.3em,topsep=2pt,itemsep=1pt,label=\textbullet]
  \item Compact metric spaces are complete (\R{3.11}).
  \item $\mathbb R^k$ is complete (\R{3.11}).
  \item $\mathbb Q$ and $(0,1)$ are \emph{not} complete (cf.\ \R{3.12}).
\end{itemize}
\end{ObjL}

\begin{ObjL}{Remark}{6.1.3}{Testing convergence without the limit}
\emph{In a complete space} --- in particular in $\mathbb R^k$ --- it is enough to check whether a
sequence is Cauchy; there is no need to know the limit in advance. In an incomplete space such as
$\mathbb Q$ this fails: a Cauchy sequence may have no limit (\xref{6.1.2}; cf.\ \R{3.11}--\R{3.12}).
\end{ObjL}

% ============================================================
\sub{6.2}{Monotone sequences}

\begin{Obj}{Definition}{6.2.1}{Monotonic sequences}
A sequence $\{s_n\}$ is monotonic in Rudin's sense (\R{3.13}) as follows:
\begin{itemize}[leftmargin=1.3em,topsep=2pt,itemsep=1pt,label=\textbullet]
  \item \emph{monotonic increasing} if $s_n\le s_{n+1}$ for all $n$;
  \item \emph{monotonic decreasing} if $s_n\ge s_{n+1}$ for all $n$.
\end{itemize}
``Monotonic'' means either increasing or decreasing.
\end{Obj}

\begin{Obj}{Theorem}{6.2.2}{Monotone convergence}
If $\{s_n\}\subseteq\mathbb R$ is monotonic and bounded, then $\{s_n\}$ converges
(\R{3.14}).
\end{Obj}
\begin{Prf}{6.2.2}{Proof by a Supremum Argument of the Monotone Convergence Theorem.}{[A5]\quad cf.\ \R{3.14}; the decreasing case follows by applying the increasing case to $\{-s_n\}$.}
\item Treat the increasing case (the decreasing case is analogous). Since $\{s_n\}$ is bounded above, the least-upper-bound property of $\mathbb R$ provides
  \[ s := \sup_n s_n . \]
\item Let $\eps>0$. Because $s-\eps<s$, the number $s-\eps$ is not an upper bound of $\{s_n\}$, so there is $N_\eps$ with
  \[ s_{N_\eps} > s-\eps . \]
\item As the sequence is increasing and $s$ is an upper bound, for every $n\ge N_\eps$
  \[ s-\eps < s_{N_\eps} \le s_n \le s , \]
  hence $|s_n-s|<\eps$. Therefore $s_n\to s$.
\end{Prf}

\begin{Fig}{6.2.3}{Climbing to the supremum}{An increasing sequence bounded above closes in on $s=\sup_n s_n$: the terms rise but never cross the dashed bound, and from some $N_\eps$ on they all lie in the band $(s-\eps,\,s]$.}
\begin{tikzpicture}[x=0.62cm,y=1.75cm]
  \fill[figLim!12] (0,1.26) rectangle (9,1.5);
  \draw[->] (0,0) -- (9.3,0) node[right] {$n$};
  \draw[->] (0,0) -- (0,1.78) node[above] {$s_n$};
  \draw[figLim,dashed,thick] (0,1.5) -- (9,1.5);
  \node[left,figLim] at (0,1.5) {$s$};
  \draw[figLim,dotted,thick] (0,1.26) -- (9,1.26);
  \node[left,figLim] at (0,1.26) {\small $s-\eps$};
  \foreach \n/\y in {1/0.5,2/0.84,3/1.06,4/1.19,5/1.31,6/1.385,7/1.43,8/1.46} {
     \filldraw[figTerm] (\n,\y) circle (1.15pt);
     \draw (\n,0.022) -- (\n,-0.022) node[below] {\scriptsize $\n$};
  }
\end{tikzpicture}
\end{Fig}

% ============================================================
\sub{6.3}{Comparison and squeeze}

\begin{Obj}{Proposition}{6.3.1}{Order is preserved in the limit}
Let $\{s_n\}$ and $\{t_n\}$ be convergent real sequences with $s_n\le t_n$ for all $n$. Then
\[ \lim_{n\to\infty}s_n \;\le\; \lim_{n\to\infty}t_n . \]
This is the ordinary-limit form of Rudin's comparison theorem for upper and lower limits
(cf.\ \R{3.19}).
\end{Obj}

\noindent\emph{Idea.} Suppose instead that $s:=\lim s_n$ and $t:=\lim t_n$ satisfy $s>t$.
Separate $s$ and $t$ by disjoint $\eps$-balls, with $\eps<\tfrac12|s-t|$, so that $\Beps(t)$
and $\Beps(s)$ sit on opposite sides of the midpoint $\tfrac{s+t}{2}$.

\begin{Fig}{6.3.2}{Separating the limits}{With $\eps<\tfrac12|s-t|$, the balls $\Beps(t)$ and $\Beps(s)$ fall on opposite sides of the midpoint $\tfrac{s+t}{2}$, which is impossible once $s_n\le t_n$.}
\begin{tikzpicture}[scale=1.1]
  \draw[->] (-0.5,0) -- (8.2,0);
  % left ball around t
  \draw[figLim,thick] (1.2,0) -- (2.8,0);
  \draw[figLim] (1.2,0.12) -- (1.2,-0.12);
  \draw[figLim] (2.8,0.12) -- (2.8,-0.12);
  \node[above,figLim] at (2,0.5) {$\Beps(t)$};
  \filldraw (2,0) circle (1.4pt) node[below=4pt] {$t$};
  % t_n is caught in the left ball, since t_n < t+eps
  \draw[figTerm] (2.5,0.1) -- (2.5,-0.1);
  \node[above=1pt,figTerm] at (2.5,0.1) {\small $t_n$};
  % midpoint
  \filldraw (4,0) circle (1pt) node[below=5pt] {$\dfrac{s+t}{2}$};
  % right ball around s
  \draw[figLim,thick] (5.2,0) -- (6.8,0);
  \draw[figLim] (5.2,0.12) -- (5.2,-0.12);
  \draw[figLim] (6.8,0.12) -- (6.8,-0.12);
  \node[above,figLim] at (6,0.5) {$\Beps(s)$};
  \filldraw (6,0) circle (1.4pt) node[below=4pt] {$s$};
  % s_n is caught in the right ball, since s_n > s-eps
  \draw[figTerm] (5.5,0.1) -- (5.5,-0.1);
  \node[above=1pt,figTerm] at (5.5,0.1) {\small $s_n$};
  % the resulting contradiction, drawn clear of the midpoint fraction
  \draw[figBad,dotted,thin] (2.5,-0.12) -- (2.5,-1.32);
  \draw[figBad,dotted,thin] (5.5,-0.12) -- (5.5,-1.32);
  \draw[figBad,<->] (2.5,-1.32) -- (5.5,-1.32);
  \node[below=1pt,figBad] at (4,-1.32) {\footnotesize forces $t_n<s_n$, contradicting $s_n\le t_n$};
\end{tikzpicture}
\end{Fig}

\begin{Prf}{6.3.1}{Proof by Contradiction that Order Passes to the Limit.}{[A9]}
\item Write $s=\lim s_n$, $t=\lim t_n$, and suppose for contradiction $s>t$. Since $s-t>0$, choose
  \[ 0 < \eps < \tfrac12|s-t| = \tfrac12(s-t). \]
\item By convergence there are $N_1,N_2$ with $|s_n-s|<\eps$ for $n\ge N_1$ and $|t_n-t|<\eps$ for $n\ge N_2$. With $N=\max\{N_1,N_2\}$, for every $n\ge N$
  \[ s_n > s-\eps \qquad\text{and}\qquad t_n < t+\eps . \]
\item The choice of $\eps$ puts these on opposite sides of the midpoint:
  \[ t+\eps < t+\tfrac12(s-t) = \frac{s+t}{2} = s-\tfrac12(s-t) < s-\eps . \]
\item Chaining the inequalities, for every $n\ge N$,
  \[ t_n < t+\eps < \frac{s+t}{2} < s-\eps < s_n , \]
  so $t_n<s_n$, contradicting $s_n\le t_n$. Hence $s\le t$, i.e.\ $\lim s_n\le\lim t_n$.
\end{Prf}

\begin{Obj}{Proposition}{6.3.3}{Squeeze theorem}
Suppose $r_n\le s_n\le t_n$ for all $n$, where $r_n\to\ell$ and $t_n\to\ell$. Then $s_n\to\ell$.
This is the comparison principle used in Rudin's special-sequence computations (cf.\ \R{3.20}).
\end{Obj}

\begin{Fig}{6.3.4}{The squeeze}{Once $r_n$ and $t_n$ lie in $(\ell-\eps,\ell+\eps)$, the trapped term $s_n$ must lie there too.}
\begin{tikzpicture}[scale=1.1]
  \fill[figLim!12] (1,-0.16) rectangle (5,0.16);
  \draw[->] (-0.5,0) -- (7.2,0);
  \node[figLim] at (1,0) {\large $($};
  \node[figLim] at (5,0) {\large $)$};
  \node[below=3pt,figLim] at (1,0) {$\ell-\eps$};
  \node[below=3pt,figLim] at (5,0) {$\ell+\eps$};
  \filldraw[figLim] (3,0) circle (1pt) node[below=3pt,figLim] {$\ell$};
  \draw[figTerm] (1.9,0.12) -- (1.9,-0.12);
  \filldraw[figTerm] (2.7,0) circle (1.3pt);
  \draw[figTerm] (3.6,0.12) -- (3.6,-0.12);
  \node[above,figTerm] at (1.9,0.12) {$r_n$};
  \node[above,figTerm] at (2.7,0.14) {$s_n$};
  \node[above,figTerm] at (3.6,0.12) {$t_n$};
\end{tikzpicture}
\end{Fig}

\begin{Prf}{6.3.3}{Proof by an $\eps$--$N$ Estimate of the Squeeze Theorem.}{[A2]}
\item Let $\eps>0$. Since $r_n\to\ell$ and $t_n\to\ell$, there is $N$ with, for all $n\ge N$,
  \[ |r_n-\ell|<\eps \qquad\text{and}\qquad |t_n-\ell|<\eps , \]
  so $r_n,t_n\in(\ell-\eps,\ell+\eps)$.
\item Using $r_n\le s_n\le t_n$, for every $n\ge N$
  \[ \ell-\eps < r_n \le s_n \le t_n < \ell+\eps , \]
  hence $|s_n-\ell|<\eps$. As $\eps>0$ was arbitrary, $s_n\to\ell$.
\end{Prf}

% ============================================================
\sub{6.4}{The upper and lower limits}

\noindent Convergence $s_n\to\ell$ says that for every $\eps>0$ there is $N$ with
$s_n\in(\ell-\eps,\ell+\eps)$ for $n\ge N$. For a bounded sequence that does \emph{not}
converge, no single $\ell$ works. Instead there are two values $\ell^-\le\ell^+$ such that the
sequence is eventually trapped in any $\eps$-enlargement of the band $[\ell^-,\ell^+]$:
\[ \forall\,\eps>0\ \ \exists\,N\ \text{ s.t. }\ s_n\in(\ell^--\eps,\ \ell^++\eps)\quad\text{for }n\ge N . \]
We write $\ell^-=\liminf s_n$ and $\ell^+=\limsup s_n$.

\begin{Fig}{6.4.1}{The trapping band}{For a bounded sequence, every $\eps$-enlargement of $[\ell^-,\ell^+]$ eventually contains all the terms.}
\begin{tikzpicture}[scale=1.1]
  \fill[figLim!10] (1,-0.16) rectangle (5.5,0.16);
  \draw[->] (-0.3,0) -- (7.2,0);
  \node[figLim] at (1,0) {\large $($};
  \node[figLim] at (5.5,0) {\large $)$};
  \node[below=3pt,figLim] at (1,0) {$\ell^- - \eps$};
  \node[below=3pt,figLim] at (5.5,0) {$\ell^+ + \eps$};
  \draw[figLim,very thick] (2,0) -- (4.5,0);
  \draw[figLim] (2,0.12) -- (2,-0.12);
  \draw[figLim] (4.5,0.12) -- (4.5,-0.12);
  \node[above=2pt,figLim] at (2,0.12) {$\ell^-$};
  \node[above=2pt,figLim] at (4.5,0.12) {$\ell^+$};
  \filldraw[figTerm] (3.2,0) circle (1.2pt);
  \node[above=2pt,figTerm] at (3.2,0.12) {$s_n$};
\end{tikzpicture}
\end{Fig}

\noindent The two coincide exactly when the sequence converges:
\[ s_n\to\ell \iff \liminf s_n = \limsup s_n = \ell . \]

\begin{ObjL}{Example}{6.4.2}{An oscillating sequence}
For $s_n=(-1)^n\dfrac{n}{n-1}$ (with $n\ge2$), the even terms approach $+1$ from above and the
odd terms approach $-1$ from below. Hence $\liminf s_n=-1$ and $\limsup s_n=+1$, and the
sequence does not converge, since these differ (cf.\ \R{3.18}).
\end{ObjL}

\begin{Fig}{6.4.3}{Two limit points}{The oscillating sequence $s_n=(-1)^n\frac{n}{n-1}$ has no limit: the even terms fall toward $\ell^+=\limsup s_n=1$ from above and the odd terms rise toward $\ell^-=\liminf s_n=-1$ from below. Every $\eps$-enlargement of the shaded band $[\ell^-,\ell^+]$ eventually traps the whole tail.}
\begin{tikzpicture}[x=0.66cm,y=0.92cm]
  \fill[figLim!12] (1,-1) rectangle (10,1);
  \draw[figLim,dashed,thick] (1,1) -- (10,1);
  \draw[figLim,dashed,thick] (1,-1) -- (10,-1);
  \node[anchor=east,font=\scriptsize,figLim] at (0.9,1) {$\ell^+\!=\!1$};
  \node[anchor=east,font=\scriptsize,figLim] at (0.9,-1) {$\ell^-\!=\!-1$};
  \draw[->] (1,0) -- (10.5,0) node[right] {$n$};
  \draw[->] (1,-1.8) -- (1,2.4) node[above] {$s_n$};
  \draw[figTerm!55,thin] (2,2) -- (3,-1.5) -- (4,1.3333) -- (5,-1.25) -- (6,1.2)
     -- (7,-1.1667) -- (8,1.1429) -- (9,-1.125);
  \foreach \n/\y in {2/2,3/-1.5,4/1.3333,5/-1.25,6/1.2,7/-1.1667,8/1.1429,9/-1.125}{
     \filldraw[figTerm] (\n,\y) circle (1.15pt);
     \draw (\n,0.06) -- (\n,-0.06) node[below=1pt] {\scriptsize $\n$};
  }
\end{tikzpicture}
\end{Fig}

\begin{Obj}{Definition}{6.4.4}{Divergence to $\pm\infty$}
For a real sequence $\{t_n\}$ (\R{3.15}),
\begin{itemize}[leftmargin=1.3em,topsep=2pt,itemsep=1pt,label=\textbullet]
  \item $t_n\to+\infty$ means: for every $M>0$ there is $N$ with $t_n>M$ for all $n\ge N$;
  \item $t_n\to-\infty$ means: for every $M>0$ there is $N$ with $t_n<-M$ for all $n\ge N$.
\end{itemize}
Such sequences diverge, but in a specific, controlled way.
\end{Obj}

\begin{Obj}{Definition}{6.4.5}{Limit inferior and limit superior}
The \emph{upper} and \emph{lower limits} of a real sequence $\{s_n\}$ are defined by the tail
extrema
\[ \limsup s_n=\lim_{n\to\infty}\Big(\sup_{k\ge n}s_k\Big),\qquad
   \liminf s_n=\lim_{n\to\infty}\Big(\inf_{k\ge n}s_k\Big). \]
The tails are nested, $\{s_k:k\ge n+1\}\subseteq\{s_k:k\ge n\}$, so $\sup_{k\ge n}s_k$ is
non-increasing in $n$ (and $\inf_{k\ge n}s_k$ non-decreasing); hence both limits always exist in
the extended reals $[-\infty,+\infty]$. Two characterizations follow.
\begin{enumerate}[label=\textbf{(\arabic*)},leftmargin=2.2em,topsep=3pt,itemsep=5pt]
  \item \textbf{Subsequential limits.} If $\{s_n\}$ is bounded and $E$ is its set of
        subsequential limits, then
        \[ \liminf s_n=\min E,\qquad \limsup s_n=\max E . \]
  \item \textbf{Eventual trapping} (finite $\liminf,\limsup$). For every $\eps>0$ there is
        $N_\eps$ with
        \[ s_n\in\big(\liminf s_n-\eps,\ \limsup s_n+\eps\big)\qquad\text{for }n\ge N_\eps , \]
and $[\liminf s_n,\ \limsup s_n]$ is the smallest closed interval with this property.
\end{enumerate}
This tail-extrema presentation is equivalent to Rudin's subsequential-limit definition
(\R{3.16}) and characterization (\R{3.17}).
\end{Obj}

\begin{ObjL}{Remark}{6.4.6}{The extreme limits are attained}
For a \emph{bounded} sequence, $E$ is closed, bounded, and nonempty, so it contains its own
extrema: $\inf E\in E$ and $\sup E\in E$. Hence $\liminf s_n$ and $\limsup s_n$ are themselves
subsequential limits of $\{s_n\}$. (For an unbounded sequence this holds in the extended reals,
with $\pm\infty$ admitted as subsequential limits.) Compare \R{3.17}.
\end{ObjL}

\begin{ObjL}{Remark}{6.4.7}{Existence}
If $\{s_n\}$ is bounded, then $E$ is bounded and nonempty, so $\liminf s_n$ and $\limsup s_n$
exist as finite reals. If $\{s_n\}$ is unbounded, then $E$ may include $+\infty$ and/or
$-\infty$, and $\liminf s_n,\limsup s_n$ may equal $\pm\infty$ (cf.\ \R{3.16}--\R{3.17}).
\end{ObjL}

% ============================================================
\sub{6.5}{Some standard limits}

\begin{Obj}{Proposition}{6.5.1}{Four standard limits}
\begin{enumerate}[label=\textup{(\alph*)},leftmargin=2em,topsep=2pt,itemsep=2pt]
  \item If $|x|<1$, then $x^n\to0$.
  \item If $p>0$, then $\dfrac{1}{n^p}\to0$.
  \item If $p>1$, then $\sqrt[n]{p}\to1$.
  \item $\sqrt[n]{n}\to1$.
\end{enumerate}
These are the recurring special limits collected in \R{3.20}(a)--(c),(e).
\end{Obj}
\begin{Prf}{6.5.1.a}{Proof by the Binomial Inequality of $x^n\to0$.}{[A2]}
\item If $x=0$ then $x^n=0\to0$, so assume $0<|x|<1$. Then $1/|x|>1$, so $|x|=\dfrac{1}{1+p}$ for a unique $p>0$ (the map $p\mapsto\frac{1}{1+p}$ is a bijection $(0,\infty)\to(0,1)$).
\item By the binomial theorem, keeping the nonnegative $k=0,1$ terms,
  \[ (1+p)^n=\sum_{k=0}^n\binom{n}{k}p^k\ \ge\ 1+np . \]
\item Hence
  \[ |x|^n=\frac{1}{(1+p)^n}\le\frac{1}{1+np}\xrightarrow[n\to\infty]{}0 , \]
  and since $|x^n|=|x|^n\to0$, we conclude $x^n\to0$.
\end{Prf}
\begin{Prf}{6.5.1.b}{Proof by an $\eps$--$N$ Estimate of $1/n^p\to0$.}{[A2]}
\item Let $\eps>0$ and choose $N$ with $N>(1/\eps)^{1/p}$, equivalently $\dfrac{1}{N^p}<\eps$.
\item Then for every $n\ge N$,
  \[ \frac{1}{n^p}\le\frac{1}{N^p}<\eps , \]
  so $\dfrac{1}{n^p}\to0$.
\end{Prf}
\begin{Prf}{6.5.1.c}{Proof by the Squeeze Theorem of $\sqrt[n]{p}\to1$.}{[A2]\quad uses \xref{6.3.3}.}
\item Assume $p>1$ and set $x_n=\sqrt[n]{p}-1\ge0$. By the binomial inequality from (a),
  \[ 1+n x_n\le(1+x_n)^n=p , \]
  so $0\le x_n\le\dfrac{p-1}{n}$.
\item Since $\dfrac{p-1}{n}\to0$, the squeeze theorem (\xref{6.3.3}) gives $x_n\to0$, i.e.\ $\sqrt[n]{p}=1+x_n\to1$.
\end{Prf}
\begin{Prf}{6.5.1.d}{Proof by the Squeeze Theorem of $\sqrt[n]{n}\to1$.}{[A2]\quad keep the $k=2$ binomial term; uses \xref{6.3.3}.}
\item Set $x_n=\sqrt[n]{n}-1\ge0$, so $n=(1+x_n)^n$. The linear term is too weak, so keep the $k=2$ term:
  \[ n=(1+x_n)^n\ge\binom{n}{2}x_n^2=\frac{n(n-1)}{2}\,x_n^2\qquad(n\ge2). \]
\item Rearranging, $x_n^2\le\dfrac{2}{n-1}$, hence $0\le x_n\le\sqrt{\dfrac{2}{n-1}}$.
\item Since $\sqrt{\tfrac{2}{n-1}}\to0$, the squeeze theorem (\xref{6.3.3}) gives $x_n\to0$, i.e.\ $\sqrt[n]{n}=1+x_n\to1$.
\end{Prf}

\begin{ObjL}{Remark}{6.5.2}{The case $0<p<1$ in (c)}
For $0<p<1$, apply part~(c) to $\sigma=1/p>1$: since $\sqrt[n]{\sigma}\to1$, also
$\sqrt[n]{p}=1/\sqrt[n]{\sigma}\to1$ (and $\sqrt[n]{1}=1$). Hence $\sqrt[n]{p}\to1$ for every
$p>0$; the full $p>0$ statement is \R{3.20}(b).
\end{ObjL}

% ============================================================
\sub{6.6}{Introduction to series}

{\itshape\small (Not covered on the midterm.)}\par\smallskip
\noindent As a field, $\mathbb R$ hands us directly only the polynomials $1,x,x^2,\dots$ and,
as their quotients, the rational functions $P(x)/Q(x)$. What about $e^x$, $\sin x$, $\cos x$?
These are reached not by field operations but by infinite sums --- for instance
\[ e=\sum_{k=0}^\infty\frac{1}{k!}. \]
To make sense of such a sum we must first say what an infinite sum \emph{is}.

\begin{Obj}{Definition}{6.6.1}{Series and partial sums}
Given a sequence $\{a_n\}$ in $\mathbb R$, its \emph{$k$-th partial sum} is
\[ s_k=a_1+\dots+a_k=\sum_{i=1}^k a_i\qquad(s_0=0,\text{ the empty sum}). \]
If $s_k\to s$, the \emph{series} $\sum a_n$ \emph{converges}, and we write
\[ \sum_{i=1}^\infty a_i=\lim_{k\to\infty}s_k=s . \]
A series may instead be indexed from any starting point---often $n=0$, as for $\sum c_n x^n$---with
the partial sums adjusted accordingly.
This is Rudin's series convention (\R{3.21}) with the real-valued specialization used in lecture.
\end{Obj}

\begin{Obj}{Theorem}{6.6.2}{Cauchy criterion for series}
A series $\sum a_n$ converges if and only if for every $\eps>0$ there is $N$ such that
\[ \Bigl|\sum_{k=m}^{n}a_k\Bigr|<\eps \qquad\text{for all } n\ge m\ge N . \]
This is the Cauchy criterion for series (\R{3.22}).
\end{Obj}
\begin{Prf}{6.6.2}{Proof by Completeness of $\mathbb R$ of the Cauchy Criterion.}{[Direct]\quad cf.\ \R{3.22}; the partial sums converge iff they are Cauchy, and $\mathbb R$ is complete.}
\item For $n\ge m\ge1$ one has $s_n-s_{m-1}=\sum_{k=m}^{n}a_k$ (with $s_0=0$), so the displayed condition says precisely that the partial sums $\{s_k\}$ form a Cauchy sequence.
\item Since $\mathbb R$ is complete (\xref{6.1.1}), $\{s_k\}$ converges if and only if it is Cauchy. Hence $\sum a_n$ converges if and only if the criterion holds.
\end{Prf}

\begin{Obj}{Corollary}{6.6.3}{Vanishing of the terms}
If $\sum a_n$ converges, then $a_n\to0$ (\R{3.23}).
\end{Obj}
\begin{Prf}{6.6.3}{Proof by the Cauchy Criterion with $m=n$.}{[Direct]\quad cf.\ \R{3.23}.}
\item Take $m=n$ in the criterion (\xref{6.6.2}): for every $\eps>0$ there is $N$ with
  \[ |a_n|=\Bigl|\sum_{k=n}^{n}a_k\Bigr|<\eps\qquad\text{for all }n\ge N . \]
\item Therefore $a_n\to0$.
\end{Prf}

\noindent So $a_n\to0$ is \emph{necessary} for convergence. Is it \emph{sufficient}? No.

\begin{Obj}{Proposition}{6.6.4}{The harmonic series diverges}
Although $\dfrac{1}{n}\to0$, the harmonic series $\displaystyle\sum_{n=1}^\infty\frac1n$ diverges.
This is the $p=1$ case of Rudin's $p$-series theorem (\R{3.28}).
\end{Obj}
\begin{Prf}{6.6.4}{Proof by Dyadic Grouping of the Divergence of the Harmonic Series.}{[Direct]\quad Cauchy condensation; cf.\ \R{3.27}--\R{3.28}.}
\item Round each term down to the next power of $\tfrac12$: replace $\tfrac13$ by $\tfrac14$; then $\tfrac15,\tfrac16,\tfrac17$ by $\tfrac18$; and so on. If $s'_k$ denotes the resulting partial sums, then $s_k\ge s'_k$ for every $k$.
\item The rounded series groups into dyadic blocks of equal weight:
  \[ 1+\tfrac12+\Bigl(\tfrac14+\tfrac14\Bigr)+\Bigl(\tfrac18+\tfrac18+\tfrac18+\tfrac18\Bigr)+\cdots
     = 1+\tfrac12+\tfrac12+\tfrac12+\cdots , \]
  each block summing to $\tfrac12$; hence $s'_{2^m}=1+\tfrac{m}{2}$.
\item Therefore $s_{2^m}\ge s'_{2^m}=1+\tfrac{m}{2}\to\infty$, so the partial sums $\{s_n\}$ have a subsequence diverging to $+\infty$ and cannot converge. Thus $\textstyle\sum\frac1n$ diverges.
\end{Prf}

\begin{Obj}{Proposition}{6.6.5}{The $p$-series}
The series $\displaystyle\sum_{n=1}^\infty\frac{1}{n^p}$ converges for $p>1$.
Rudin proves the full convergence/divergence dichotomy in \R{3.28}.
\end{Obj}
\PrfRef{6.6.5}{Proof of the Convergence of the $p$-Series.}{Stated in lecture as a fact, for use as a comparison series; cf.\ \R{3.28}.}

\begin{Obj}{Proposition}{6.6.6}{A nonnegative series converges when bounded}
A series $\sum a_n$ of nonnegative terms ($a_n\ge0$) converges if and only if its partial sums
are bounded (\R{3.24}).
\end{Obj}
\begin{Prf}{6.6.6}{Proof by Monotone Convergence for a Series of Nonnegative Terms.}{[A5]\quad cf.\ \R{3.24}; the partial sums increase; apply \xref{6.2.2}.}
\item Since $a_{n+1}\ge0$, the partial sums satisfy $s_{n+1}=s_n+a_{n+1}\ge s_n$, so $\{s_n\}$ is increasing.
\item If the partial sums are bounded, then $\{s_n\}$ is increasing and bounded, hence converges by the monotone convergence theorem (\xref{6.2.2}); that is, $\sum a_n$ converges. Conversely, a convergent sequence is bounded.
\end{Prf}

\begin{Obj}{Proposition}{6.6.7}{The geometric series}
For $|x|<1$,
\[ \sum_{n=0}^\infty x^n = \frac{1}{1-x} . \]
This is the convergent half of Rudin's geometric-series theorem (\R{3.26}).
\end{Obj}
\begin{Prf}{6.6.7}{Proof by the Partial-Sum Formula of the Geometric Series.}{[Direct]\quad cf.\ \R{3.26}; uses $|x|^{k}\to0$ for $|x|<1$ (\xref{6.5.1}).}
\item For $x\ne1$ the partial sums have the closed form
  \[ s_k = 1+x+\dots+x^k = \frac{1-x^{k+1}}{1-x} . \]
\item Since $|x|<1$ gives $|x|^{k+1}\to0$, we have $x^{k+1}\to0$ and hence
  \[ s_k \xrightarrow[k\to\infty]{} \frac{1}{1-x} , \]
  so $\displaystyle\sum_{n=0}^\infty x^n = \frac{1}{1-x}$.
\end{Prf}

\begin{Obj}{Theorem}{6.6.8}{Comparison test for series}
\begin{enumerate}[label=\textup{(\alph*)},leftmargin=2em,topsep=2pt,itemsep=2pt]
  \item If $|a_n|\le c_n$ for all $n\ge N_0$ and $\sum c_n$ converges, then $\sum a_n$ converges.
  \item If $a_n\ge d_n\ge0$ for all $n$ and $\sum d_n$ diverges, then $\sum a_n$ diverges.
\end{enumerate}
This is Rudin's comparison test (\R{3.25}).
\end{Obj}
\begin{Prf}{6.6.8}{Proof by the Cauchy Criterion of the Comparison Test.}{[A4]\quad cf.\ \R{3.25}; bound the tail of $\sum a_n$ by that of $\sum c_n$.}
\item For (a), let $\eps>0$. Since $\sum c_n$ converges, the Cauchy criterion (\xref{6.6.2}) gives some $N\ge N_0$ with $\sum_{k=m}^n c_k<\eps$ for all $n\ge m\ge N$. Then for such $n,m$,
  \[ \Bigl|\sum_{k=m}^n a_k\Bigr| \le \sum_{k=m}^n |a_k| \le \sum_{k=m}^n c_k < \eps , \]
  so $\sum a_n$ meets the Cauchy criterion and therefore converges.
\item For (b), suppose instead $\sum a_n$ converged. Since $0\le d_n\le a_n=|a_n|$, part~(a) (with $c_n=a_n$) would force $\sum d_n$ to converge --- contradicting that it diverges. Hence $\sum a_n$ diverges.
\end{Prf}

\noindent The natural series to compare against is the geometric one. Comparing $\sum a_n$ to a
geometric $\sum b_n$: if the ratio $b_{n+1}/b_n$ is constant we are led to the \emph{ratio test},
and if the root $\sqrt[n]{b_n}$ is constant, to the \emph{root test}.

\begin{Obj}{Theorem}{6.6.9}{Root test}
Given $\sum a_n$, set $\alpha=\limsup_{n\to\infty}\sqrt[n]{|a_n|}$. Then
\begin{itemize}[leftmargin=1.3em,topsep=2pt,itemsep=1pt,label=\textbullet]
  \item if $\alpha<1$, $\sum a_n$ converges;
  \item if $\alpha>1$, $\sum a_n$ diverges;
  \item if $\alpha=1$, the test is inconclusive.
\end{itemize}
This is Rudin's root test (\R{3.33}).
\end{Obj}
\begin{Prf}{6.6.9}{Proof by Comparison with a Geometric Series of the Root Test.}{[A1]\quad cf.\ \R{3.33}; uses the comparison test (\xref{6.6.8}).}
\item If $\alpha<1$, choose $\beta$ with $\alpha<\beta<1$. Then $\sqrt[n]{|a_n|}<\beta$ for all large $n$, so $|a_n|<\beta^n$; as $\sum\beta^n$ converges, the comparison test (\xref{6.6.8}) gives convergence.
\item If $\alpha>1$, then $\sqrt[n]{|a_n|}>1$ for infinitely many $n$, so $|a_n|>1$ infinitely often and $a_n\not\to0$; by \xref{6.6.3} the series diverges. For $\alpha=1$ the test decides nothing: both $\sum\frac1n$ (divergent) and $\sum\frac1{n^2}$ (convergent) have $\alpha=1$.
\end{Prf}

\begin{Obj}{Theorem}{6.6.10}{Ratio test}
Suppose $a_n\neq0$ for all large $n$. Then $\sum a_n$ converges if
$\displaystyle\limsup_{n\to\infty}\Bigl|\frac{a_{n+1}}{a_n}\Bigr|<1$, and diverges if there is
$N$ with $\Bigl|\dfrac{a_{n+1}}{a_n}\Bigr|\ge1$ for all $n\ge N$.
This is Rudin's ratio test (\R{3.34}).
\end{Obj}
\begin{Prf}{6.6.10}{Proof by Comparison with a Geometric Series of the Ratio Test.}{[A1]\quad cf.\ \R{3.34}; uses the comparison test (\xref{6.6.8}).}
\item If $\limsup|a_{n+1}/a_n|<1$, choose $\beta$ in between; then $|a_{n+1}/a_n|<\beta$ for $n\ge N$, so $|a_n|\le|a_N|\,\beta^{\,n-N}$ for $n\ge N$. Comparison with the geometric series (\xref{6.6.8}) gives convergence.
\item If $|a_{n+1}/a_n|\ge1$ for all $n\ge N$, then $|a_n|\ge|a_N|>0$ for $n\ge N$, so $a_n\not\to0$ and the series diverges (\xref{6.6.3}).
\end{Prf}

\begin{Obj}{Definition}{6.6.11}{Power series and radius of convergence}
A \emph{power series} is a function defined by
\[ f(x)=\sum_{n=0}^\infty c_n x^n = c_0+c_1x+c_2x^2+c_3x^3+\cdots . \]
Fix $x\neq0$ and apply the root test (\xref{6.6.9})\note{The root test is used here because it is
strictly more general than the ratio test; for the precise relationship see the appendix,
\appref{6.A.1}.} to $\sum c_n x^n$. With
$\alpha:=\limsup_{n\to\infty}\sqrt[n]{|c_n|}$,
\[ \beta:=\limsup_{n\to\infty}\sqrt[n]{|c_n x^n|}=|x|\,\alpha , \]
so the series converges when $\beta=|x|\,\alpha<1$ and diverges when $|x|\,\alpha>1$. Hence it
converges for $|x|<R$, where
\[ R:=\frac{1}{\alpha}\qquad\bigl(R:=+\infty\ \text{if}\ \alpha=0,\qquad R:=0\ \text{if}\ \alpha=+\infty\bigr). \]
$R$ is the \emph{radius of convergence}: the series converges on $(-R,R)$ and diverges for
$|x|>R$; at $x=0$ it is just $f(0)=c_0$, and the endpoints $|x|=R$ require separate analysis.
This is Rudin's power-series setup and radius formula (\R{3.38}--\R{3.39}), restricted here
to real $x$.
\end{Obj}

\begin{ObjL}{Example}{6.6.12}{The exponential series}
For $f(x)=\displaystyle\sum_{n=0}^\infty\frac{x^n}{n!}$ the ratios have a limit, so they compute
$\alpha$ (the valid case for the ratio shortcut):
\[ \lim_{n\to\infty}\frac{1/(n+1)!}{1/n!}=\lim_{n\to\infty}\frac{1}{n+1}=0=\alpha , \]
so $R=1/\alpha=+\infty$: the series converges for every $x$. This function is $f(x)=e^x$.
Compare Rudin's examples following the radius-of-convergence theorem (\R{3.40}).
\end{ObjL}

% ============================================================
\lectureappendix

\begin{ObjL}{Supplement}{6.A.1}{The root test subsumes the ratio test}
For any sequence $\{c_n\}$ with $c_n\neq0$,
\[ \liminf_{n}\Bigl|\tfrac{c_{n+1}}{c_n}\Bigr|\ \le\ \liminf_n\sqrt[n]{|c_n|}\ \le\
   \limsup_n\sqrt[n]{|c_n|}\ \le\ \limsup_n\Bigl|\tfrac{c_{n+1}}{c_n}\Bigr| \]
(\R{3.37}). Hence whenever the ratio test settles convergence the root test settles it too, but
not conversely; and \emph{if} $\lim_n|c_{n+1}/c_n|=L$ exists then the root quantity tends to $L$
as well, so the ratio test gives the same radius $R=1/L$ (the converse fails --- the radii can
agree without the ratio limit existing). This is why \xref{6.6.11} fixes $R$ through the root
quantity $\alpha=\limsup_n\sqrt[n]{|c_n|}$.
\end{ObjL}

% per-lecture Index of Objects -- only when compiling this lecture on its own;
% in the book the master index in main.tex covers all lectures.
\ifSubfilesClassLoaded{\clearpage\objindex}{}

\end{document}
